\section{Discussion and limitations}
The S1 observation distinguishes accurate correction from useful initialization. Convergence at the specified residual shows that the combined solver reaches an accurate discrete solution; it does not attribute that result to learning. At the strict endpoint, all observed methods exhaust the 63-dimensional iteration count, leaving no measured iteration advantage. This is consistent with small initial errors retaining a broad spectral support, but an observed count alone does not prove a unique spectral explanation.

Source linearity is a principled restriction that can be proved independently of experiment. Its computational value must still survive tests: a better predictor may not reduce accepted iterations, a reduction may not cover inference and setup costs, and an iterative path may remain slower than a direct method. The classical Laplacian condition-number bound also cautions against attributing gains caused by a preconditioner to the fitted predictor.

The study is deliberately narrow. Smooth four-coordinate coefficients, one-dimensional domains, modest grids and synthetic right-hand sides do not represent general PDE simulation. The forcing-shift test is limited to declared added modes. The regression model is not a deep neural operator. Five confirmation training blocks provide limited uncertainty resolution; randomized repeated timings remain host-specific. Offline training and model-loading costs are not amortized into a hypothetical workload. We make no state-of-the-art claim and no claim of independent replication by another researcher.

S1's pooled spatial calibration scores do not establish distribution-free conformal coverage for independent PDE instances. Those values remain in the archived evidence, but the coverage guarantee is withdrawn and the mechanism is omitted from S2. The source-linear predictor supplies a point estimate; certified numerical residual acceptance is a different property from probabilistic uncertainty calibration.

\section{Reproducibility and research integrity}
The repository preserves S1, a flaw register, source and protocol hashes, all admitted/failed attempts, raw synthetic inputs, trained coefficients, solver histories, timing passes and analysis. Confirmation is not used for subsequent tuning. A valid negative result closes the stated hypothesis; any new mechanism requires a successor version. Code replay and external scientific reproduction are reported separately. No human-subject data, personally identifying training data or proprietary external dataset is used.

\section{Conclusion}
The predecessor demonstrates that a convergent learned hybrid solver can offer no iteration savings. S2 tests a precise structural repair against strong numerical controls under a prospective protocol. At this pre-result stage its outcome remains unknown; the final conclusion must be determined by verified measurements, not the preference for a positive result.
